%% 第1巻 巻末（こたえ・きごうカード・たんいのひょう・チェックひょう・さくいん・奥付）
\makeatletter
%% ---------------- こたえ ----------------
\bkAnswerPages

%% ---------------- きごうカード ----------------
\bkKigoPagesOne

%% ---------------- たんいのひょう ----------------
\bkbacktitle{たんいのひょう}
{\small この \R{本}{ほん}で でてきた たんいと、 その \R{大}{おお}きさの めやすです。}\par\vspace{2mm}

\def\bkuhead#1#2{\par\addvspace{3mm}\noindent{\sffamily\bfseries\large\color{bkTheme}#1}\hspace{0.6\zw}{\small\color{bkGray}#2}\par\nobreak\vspace{1mm}}
\newcommand\bkutable[1]{\begin{bkcenter}\small
\begin{tblr}{colspec={Q[c,m,wd=22mm] Q[l,m,wd=46mm] Q[l,m,wd=58mm]}, row{1}={bg=bkVol!14, font=\sffamily\bfseries\footnotesize},
  hlines={0.4pt,bkLine}, rowsep=2pt}
たんい & よみかた & めやす・ くらべかた \\
#1
\end{tblr}\end{bkcenter}}

\bkuhead{ながさ}{（1-8）}
\bkutable{%
1\,mm & ミリメートル & 1\,cm を 10に \R{分}{わ}けた 1つ \\
1\,cm & センチメートル & ゆびの はばほど。 1\,cm ＝ 10\,mm \\
1\,m & メートル & りょう\R{手}{て}を ひろげた ほど。 1\,m ＝ 100\,cm \\}

\bkuhead{かさ}{（1-9）}
\bkutable{%
1\,mL & ミリリットル & \R{小}{こ}さじ 1ぱいは 5\,mL \\
1\,dL & デシリットル & コップ \R{半分}{はん|ぶん}ほど。 1\,dL ＝ 100\,mL \\
1\,L & リットル & \R{牛乳}{ぎゅう|にゅう}パック 1\R{本}{ほん}。 1\,L ＝ 10\,dL ＝ 1000\,mL \\}

\bkuhead{おもさ}{（1-9）}
\bkutable{%
1\,g & グラム & 1\R{円}{えん}\R{玉}{だま} 1まい \\
1\,kg & キログラム & \R{水}{みず} 1\,L。 1\,kg ＝ 1000\,g \\}

\clearpage
\bkuhead{\R{時間}{じ|かん}}{（1-10）}
\bkutable{%
1\R{秒}{びょう} & びょう & 「いち」と いう あいだの \R{時間}{じ|かん} \\
1\R{分}{ぷん} & ふん・ぷん & 1\R{分}{ぷん} ＝ 60\R{秒}{びょう} \\
1\R{時間}{じ|かん} & じかん & 1\R{時間}{じ|かん} ＝ 60\R{分}{ぷん} \\
1\R{日}{にち} & にち & 1\R{日}{にち} ＝ 24\R{時間}{じ|かん} \\}

\bkuhead{あたたかさ（\R{温度}{おん|ど}）}{（1-11）}
\bkutable{%
1\,℃ & ど & \R{水}{みず}が こおるのは 0\,℃、 にえるのは 100\,℃ \\}

\bkuhead{\R{角}{かく}の \R{大}{おお}きさ}{（1-20）}
\bkutable{%
1° & ど & 1\R{回転}{かい|てん}を 360に \R{分}{わ}けた 1つ。 \R{直角}{ちょっ|かく} ＝ 90°、 1\R{回転}{かい|てん} ＝ 360° \\}

\bkuhead{\R{広}{ひろ}さ（\R{面積}{めん|せき}）}{（1-21）}
\bkutable{%
1\,cm² & へいほう センチメートル & 1\,cm の \R{正方形}{せい|ほう|けい}の \R{広}{ひろ}さ \\
1\,m² & へいほう メートル & 1\,m の \R{正方形}{せい|ほう|けい}の \R{広}{ひろ}さ \\}

\clearpage
\bkuhead{\R{九九}{く|く}の \R{表}{ひょう}}{（1-13）}
\begin{bkcenter}
\begin{tikzpicture}[x=1mm, y=1mm]
  \foreach \a in {1,...,9} {
    \fill[bkVol!18] (\a*13,-1*11) rectangle ++(13,11);
    \node[font=\sffamily\bfseries\small, text=bkDeep] at (\a*13+6.5,-5.5) {\a};
    \fill[bkVol!18] (0,-\a*11-11) rectangle ++(13,11);
    \node[font=\sffamily\bfseries\small, text=bkDeep] at (6.5,-\a*11-5.5) {\a};
    \foreach \b in {1,...,9} {\pgfmathtruncatemacro\pr{\a*\b}
      \ifnum\a=\b \fill[bkYellow!25] (\b*13,-\a*11-11) rectangle ++(13,11);\fi
      \node[font=\small, text=bkInk] at (\b*13+6.5,-\a*11-5.5) {\pr};}}
  \fill[bkVol!35] (0,-11) rectangle ++(13,11);
  \node[font=\sffamily\bfseries\footnotesize, text=bkDeep] at (6.5,-5.5) {×};
  \draw[bkLine, line width=0.4pt] (0,0) grid[xstep=13, ystep=11] (130,-110);
  \draw[bkVol, line width=1pt] (0,0) rectangle (130,-110) (13,0) -- (13,-110) (0,-11) -- (130,-11);
\end{tikzpicture}
\end{bkcenter}
{\small \R{左}{ひだり}の たての \R{数}{かず}が 「1つ\R{分}{ぶん}」、 \R{上}{うえ}の よこの \R{数}{かず}が 「いくつ\R{分}{ぶん}」です。 きいろの ななめの ますは、 \R{同}{おな}じ \R{数}{かず}どうしの かけ\R{算}{ざん}です。}

\vspace{4mm}
\bkuhead{\R{九九}{く|く}の きまり}{（1-13）}
\begin{bkcenter}
\begin{tikzpicture}[x=1mm, y=1mm]
  \foreach \y/\c/\t/\b/\e in {
    0/bkRoute2/{ふえかたの\\きまり}/{いくつ\R{分}{ぶん}が 1 ふえると、 1つ\R{分}{ぶん}だけ ふえる。}/{3 × 4 ＝ 12\enspace → \enspace 3 × 5 ＝ 15},
    -17/bkRoute4/{\R{入}{い}れかえても\\\R{同}{おな}じ}/{1つ\R{分}{ぶん}と いくつ\R{分}{ぶん}を \R{入}{い}れかえても \R{同}{おな}じ。}/{3 × 4 ＝ 4 × 3},
    -34/bkRoute5/{ななめの\\きまり}/{\R{同}{おな}じ かずどうしの かけ\R{算}{ざん}が ならぶ。}/{1 × 1、\enspace 2 × 2、\enspace 3 × 3、\enspace …}} {
    \fill[\c!12, rounded corners=2mm] (0,\y) rectangle ++(130,-15);
    \fill[\c, rounded corners=2mm] (0,\y) rectangle ++(3,-15);
    \node[anchor=west, font=\sffamily\bfseries\small, text=\c!80!black, inner sep=0pt, align=left] at (7,\y-7.5) {\t};
    \node[anchor=west, font=\small, text=bkInk, text width=93mm, inner sep=0pt, align=left] at (35,\y-7.5) {\b\\\sffamily\bfseries \e};}
\end{tikzpicture}
\end{bkcenter}

%% ---- 第5回で追加：かずの くらいと ひっ算（1-2・1-7） ----
\clearpage
\bkuhead{かずの くらい}{（1-2・1-7）}
\begin{bkcenter}
\begin{tikzpicture}[x=1mm, y=1mm]
  \foreach \i/\n/\c in {0/{\R{一万}{いち|まん}}/bkRoute5,1/{\R{千}{せん}}/bkGreen,2/{\R{百}{ひゃく}}/bkRoute4,3/{\R{十}{じゅう}}/bkBlue,4/{\R{一}{いち}}/bkOrange} {
    \fill[\c!14, rounded corners=1.5mm] (\i*25,0) rectangle ++(23,-30);
    \node[font=\sffamily\bfseries\small, text=\c!80!black] at (\i*25+11.5,-6) {\n};}
  \node[font=\fontsize{20}{20}\selectfont] at (11.5,-19) {10000};
  \node[font=\fontsize{20}{20}\selectfont] at (36.5,-19) {1000};
  \node[font=\fontsize{20}{20}\selectfont] at (61.5,-19) {100};
  \node[font=\fontsize{20}{20}\selectfont] at (86.5,-19) {10};
  \node[font=\fontsize{20}{20}\selectfont] at (111.5,-19) {1};
  \foreach \i in {0,...,3} {\draw[bk arrow, bkGray, line width=0.9pt] (\i*25+46,-33) to[bend left=35] node[below, font=\footnotesize, text=bkInk, inner sep=1.5pt] {10こで} (\i*25+21,-33);}
\end{tikzpicture}
\end{bkcenter}
{\small \R{上}{うえ}の \R{字}{じ}は くらいの \R{名前}{な|まえ}です（\R{一}{いち}の くらい、 \R{十}{じゅう}の くらい …）。 1が 10こで 10、 10が 10こで 100。 1つ \R{左}{ひだり}の くらいは、 10こ\R{分}{ぶん}です。}\par

\bkuhead{ひっ\R{算}{ざん}の しかた}{（1-7）}
\begin{bkcenter}
\begin{tikzpicture}[x=1mm, y=1mm, dg/.style={font=\fontsize{22}{22}\selectfont, text=bkInk, inner sep=0pt, anchor=base}]
  % たし算 36 ＋ 27
  \begin{scope}[shift={(19,0)}]
    \fill[bkBlue!10] (8,-31) rectangle ++(11,33); \fill[bkOrange!10] (19,-31) rectangle ++(11,33);
    \node[font=\sffamily\bfseries\scriptsize, text=bkBlue!80!black] at (13.5,4.5) {\R{十}{じゅう}};
    \node[font=\sffamily\bfseries\scriptsize, text=bkOrange!85!black] at (24.5,4.5) {\R{一}{いち}};
    \node[dg] at (13.5,-8) {3}; \node[dg] at (24.5,-8) {6};
    \node[dg] at (3,-19) {＋}; \node[dg] at (13.5,-19) {2}; \node[dg] at (24.5,-19) {7};
    \draw[bkInk, line width=1.1pt] (0,-22) -- (30,-22);
    \node[dg] at (13.5,-30) {6}; \node[dg] at (24.5,-30) {3};
    \node[font=\fontsize{11}{11}\selectfont, text=bkRed, inner sep=0pt] at (9.6,-12.8) {1};
  \end{scope}
  \node[anchor=north west, text width=64mm, font=\footnotesize, align=left] at (0,-36) {① \R{一}{いち}の くらい：6 ＋ 7 ＝ 13\\ ② 10を \R{十}{じゅう}の くらいへ（くり\R{上}{あ}がり）\\ ③ \R{十}{じゅう}の くらい：1 ＋ 3 ＋ 2 ＝ 6};
  % ひき算 52 － 27
  \begin{scope}[shift={(91,0)}]
    \fill[bkBlue!10] (8,-31) rectangle ++(11,33); \fill[bkOrange!10] (19,-31) rectangle ++(11,33);
    \node[font=\sffamily\bfseries\scriptsize, text=bkBlue!80!black] at (13.5,4.5) {\R{十}{じゅう}};
    \node[font=\sffamily\bfseries\scriptsize, text=bkOrange!85!black] at (24.5,4.5) {\R{一}{いち}};
    \node[dg] at (13.5,-8) {5}; \node[dg] at (24.5,-8) {2};
    \draw[bkRed, line width=0.9pt] (10.8,-9) -- (16.2,-1.2);
    \node[font=\fontsize{11}{11}\selectfont, text=bkRed, inner sep=0pt] at (17.6,-0.6) {4};
    \node[dg] at (3,-19) {－}; \node[dg] at (13.5,-19) {2}; \node[dg] at (24.5,-19) {7};
    \draw[bkInk, line width=1.1pt] (0,-22) -- (30,-22);
    \node[dg] at (13.5,-30) {2}; \node[dg] at (24.5,-30) {5};
  \end{scope}
  \node[anchor=north west, text width=64mm, font=\footnotesize, align=left] at (72,-36) {① 2から 7は ひけない\\ ② \R{十}{じゅう}の くらいから 10を 1つ もらう（くり\R{下}{さ}がり）。 12 － 7 ＝ 5\\ ③ \R{十}{じゅう}の くらい：4 － 2 ＝ 2};
\end{tikzpicture}
\end{bkcenter}
{\small くらいを たてに そろえて、 \R{一}{いち}の くらいから けいさんします。 こたえは、 たし\R{算}{ざん}と ひき\R{算}{ざん}で たしかめよう（36 ＋ 27 ＝ 63 なら、 63 － 27 ＝ 36）。}

%% ---------------- チェックひょう ----------------
\bkbacktitle{チェックひょう}
{\small \R{章}{しょう}の おわりの 「ふりかえり」を あつめた \R{表}{ひょう}です。 できるように なったら □に チェックを つけよう。 できなかったら、 その ページに もどって よみなおそう。}\par\vspace{2mm}
%% 自動生成（tools/mkchecklist.py）。各章の「ふりかえり」の項目を集めたもの
\begingroup\footnotesize
\def\bkcb{\tikz[baseline=-0.5ex]\draw[bkTheme, line width=0.7pt, rounded corners=0.5pt] (0,-0.3em) rectangle (0.8em,0.5em);}
\begin{longtblr}{colspec={Q[c,m,wd=9mm] X[l,m] Q[c,m,wd=10mm] Q[c,m,wd=9mm] Q[c,m,wd=9mm]}, width=\linewidth, rowhead=1, colsep=3pt,
  row{1}={bg=bkVol!14, font=\sffamily\bfseries\scriptsize}, hlines={0.4pt,bkLine}, rowsep=0.7pt}
\R{章}{しょう} & できるように なった こと & ページ & 1\R{回}{かい}め & 2\R{回}{かい}め \\
{\sffamily\bfseries\color{bkTheme}1-1} & 1つずつ ペアに して くらべられる。 & \pageref{bk:1-1:pea} & \bkcb & \bkcb \\
 & かぞえて、 すうじで かける。 & \pageref{bk:1-1:kazoeru} & \bkcb & \bkcb \\
 & おなじ かずは、 おなじ すうじ。 & \pageref{bk:1-1:onaji} & \bkcb & \bkcb \\
{\sffamily\bfseries\color{bkTheme}1-2} & 10ずつ たばねて かぞえられる。 & \pageref{bk:1-2:tabaneru} & \bkcb & \bkcb \\
 & 23の 2と 3の いみが わかる。 & \pageref{bk:1-2:kurai} & \bkcb & \bkcb \\
 & 100、 1000、 10000が わかる。 & \pageref{bk:1-2:hyaku} & \bkcb & \bkcb \\
{\sffamily\bfseries\color{bkTheme}1-3} & なにも ない ことを 0で あらわせる。 & \pageref{bk:1-3:nai} & \bkcb & \bkcb \\
 & 0に あわせて ながさを よめる。 & \pageref{bk:1-3:awaseru} & \bkcb & \bkcb \\
 & なにも ない くらいに 0を かける。 & \pageref{bk:1-3:kurai} & \bkcb & \bkcb \\
{\sffamily\bfseries\color{bkTheme}1-4} & 「あわせる」「ふえる」を しきに かける。 & \pageref{bk:1-4:awaseru} & \bkcb & \bkcb \\
 & ものさしの うえで たし\R{算}{ざん}が できる。 & \pageref{bk:1-4:monosashi} & \bkcb & \bkcb \\
 & 10を こえる たし\R{算}{ざん}が できる。 & \pageref{bk:1-4:kuriagari} & \bkcb & \bkcb \\
{\sffamily\bfseries\color{bkTheme}1-5} & 「のこり」「ちがい」を しきに かける。 & \pageref{bk:1-5:nokori} & \bkcb & \bkcb \\
 & どれだけ かわったかを もとめられる。 & \pageref{bk:1-5:kawatta} & \bkcb & \bkcb \\
 & 10の まとまりを くずして ひける。 & \pageref{bk:1-5:kuzusu} & \bkcb & \bkcb \\
{\sffamily\bfseries\color{bkTheme}1-6} & \R{数直線}{すう|ちょく|せん}で たし\R{算}{ざん}・ひき\R{算}{ざん}が できる。 & \pageref{bk:1-6:susumu} & \bkcb & \bkcb \\
 & 2つの かずの ちがいを みつけられる。 & \pageref{bk:1-6:aida} & \bkcb & \bkcb \\
 & 1\R{目}{め}もりが 10や 100でも よめる。 & \pageref{bk:1-6:ookii} & \bkcb & \bkcb \\
{\sffamily\bfseries\color{bkTheme}1-7} & \R{何十}{なん|じゅう}・\R{何百}{なん|びゃく}の けいさんが できる。 & \pageref{bk:1-7:matomari} & \bkcb & \bkcb \\
 & くらいを そろえて、 ひっ\R{算}{さん}が できる。 & \pageref{bk:1-7:kuraigoto} & \bkcb & \bkcb \\
 & くり\R{上}{あ}がり・くり\R{下}{さ}がりが わかる。 & \pageref{bk:1-7:kuriagari} & \bkcb & \bkcb \\
{\sffamily\bfseries\color{bkTheme}1-8} & 「1\,cm の いくつ\R{分}{ぶん}」で ながさを いえる。 & \pageref{bk:1-8:cm} & \bkcb & \bkcb \\
 & ものさしで cm と mm を よめる。 & \pageref{bk:1-8:yomu} & \bkcb & \bkcb \\
 & ながさの たし\R{算}{ざん}・ひき\R{算}{ざん}が できる。 & \pageref{bk:1-8:keisan} & \bkcb & \bkcb \\
{\sffamily\bfseries\color{bkTheme}1-9} & \R{水}{みず}の かさを、 L・dL・mL で いえる。 & \pageref{bk:1-9:kasa} & \bkcb & \bkcb \\
 & てんびんで、 おもさを くらべられる。 & \pageref{bk:1-9:tenbin} & \bkcb & \bkcb \\
 & おもさを、 g・kg で いえる。 & \pageref{bk:1-9:g} & \bkcb & \bkcb \\
{\sffamily\bfseries\color{bkTheme}1-10} & \mbox{\R{時計}{と|けい}}を よんで、 \R{時}{じ}こくが いえる。 & \pageref{bk:1-10:shikumi} & \bkcb & \bkcb \\
 & \R{時}{じ}こくと \mbox{\R{時間}{じ|かん}}の ちがいが わかる。 & \pageref{bk:1-10:jikan} & \bkcb & \bkcb \\
 & 60で くり\R{上}{あ}がる けいさんが できる。 & \pageref{bk:1-10:rokujuu} & \bkcb & \bkcb \\
{\sffamily\bfseries\color{bkTheme}1-11} & \mbox{1\R{目}{め}もり}が いくつかを \R{見}{み}つけられる。 & \pageref{bk:1-11:hitomemori} & \bkcb & \bkcb \\
 & \mbox{\R{温度計}{おん|ど|けい}}や はかりの \R{目}{め}もりを よめる。 & \pageref{bk:1-11:ondokei} & \bkcb & \bkcb \\
 & \R{目}{め}もりの あいだを 「だいたい」で よめる。 & \pageref{bk:1-11:daitai} & \bkcb & \bkcb \\
{\sffamily\bfseries\color{bkTheme}1-12} & 「1つ\R{分}{ぶん} × いくつ\R{分}{ぶん}」の しきが かける。 & \pageref{bk:1-12:matomari} & \bkcb & \bkcb \\
 & かけ\R{算}{ざん}を おはじきの \R{四角}{し|かく}で あらわせる。 & \pageref{bk:1-12:shikaku} & \bkcb & \bkcb \\
 & 2のだん〜5のだんの \R{九九}{く|く}を つかえる。 & \pageref{bk:1-12:dan} & \bkcb & \bkcb \\
{\sffamily\bfseries\color{bkTheme}1-13} & \R{九九}{く|く}の \R{表}{ひょう}で、 こたえを \R{見}{み}つけられる。 & \pageref{bk:1-13:yomikata} & \bkcb & \bkcb \\
 & \R{表}{ひょう}の きまりを \R{見}{み}つけられる。 & \pageref{bk:1-13:fuekata} & \bkcb & \bkcb \\
 & \R{入}{い}れかえても、 こたえは \R{同}{おな}じと わかる。 & \pageref{bk:1-13:irekae} & \bkcb & \bkcb \\
{\sffamily\bfseries\color{bkTheme}1-14} & \R{分}{わ}ける ことを、 わり\R{算}{ざん}の しきに かける。 & \pageref{bk:1-14:hitoribun} & \bkcb & \bkcb \\
 & わり\R{算}{ざん}の こたえを、 \R{九九}{く|く}で \R{見}{み}つける。 & \pageref{bk:1-14:kuku} & \bkcb & \bkcb \\
 & 0で わる わり\R{算}{ざん}は できないと わかる。 & \pageref{bk:1-14:zero} & \bkcb & \bkcb \\
{\sffamily\bfseries\color{bkTheme}1-15} & あまりの ある わり\R{算}{ざん}が できる。 & \pageref{bk:1-15:amari} & \bkcb & \bkcb \\
 & わり\R{算}{ざん}の こたえを たしかめられる。 & \pageref{bk:1-15:tashikame} & \bkcb & \bkcb \\
 & くり\R{返}{かえ}す ものを わり\R{算}{ざん}で しらべられる。 & \pageref{bk:1-15:shuu} & \bkcb & \bkcb \\
{\sffamily\bfseries\color{bkTheme}1-16} & はんぱを \mbox{\R{小数}{しょう|すう}}で あらわせる。 & \pageref{bk:1-16:shousuu} & \bkcb & \bkcb \\
 & \mbox{\R{数直線}{すう|ちょく|せん}}で \mbox{\R{小数}{しょう|すう}}が よめる。 & \pageref{bk:1-16:suuchokusen} & \bkcb & \bkcb \\
 & \mbox{\R{半分}{はん|ぶん}}を \mbox{\R{分数}{ぶん|すう}}で あらわせる。 & \pageref{bk:1-16:bunsuu} & \bkcb & \bkcb \\
{\sffamily\bfseries\color{bkTheme}1-17} & しらべて かぞえ、 \R{表}{ひょう}に まとめられる。 & \pageref{bk:1-17:hyou} & \bkcb & \bkcb \\
 & ぼうグラフが かけて、 よめる。 & \pageref{bk:1-17:bou} & \bkcb & \bkcb \\
 & グラフの 1\R{目}{め}もりを たしかめられる。 & \pageref{bk:1-17:memori} & \bkcb & \bkcb \\
{\sffamily\bfseries\color{bkTheme}1-18} & \R{表}{ひょう}を、 おれ\R{線}{せん}グラフに かける。 & \pageref{bk:1-18:kaku} & \bkcb & \bkcb \\
 & \R{線}{せん}の むきで、 ふえる・へるが わかる。 & \pageref{bk:1-18:muki} & \bkcb & \bkcb \\
 & かたむきで、 かわる はやさが わかる。 & \pageref{bk:1-18:kyuu} & \bkcb & \bkcb \\
{\sffamily\bfseries\color{bkTheme}1-19} & \R{三角形}{さん|かく|けい}と \R{四角形}{し|かく|けい}を \R{見}{み}わけられる。 & \pageref{bk:1-19:sankaku} & \bkcb & \bkcb \\
 & \R{直角}{ちょっ|かく}や \R{長方形}{ちょう|ほう|けい}・\R{正方形}{せい|ほう|けい}が わかる。 & \pageref{bk:1-19:chokkaku} & \bkcb & \bkcb \\
 & \R{円}{えん}の \R{中心}{ちゅう|しん}と \R{半径}{はん|けい}が わかる。 & \pageref{bk:1-19:en} & \bkcb & \bkcb \\
{\sffamily\bfseries\color{bkTheme}1-20} & \R{角}{かく}は 「まわった \R{大}{おお}きさ」と わかる。 & \pageref{bk:1-20:kaku} & \bkcb & \bkcb \\
 & 1\R{回転}{かい|てん}は 360°、 \R{直角}{ちょっ|かく}は 90°と わかる。 & \pageref{bk:1-20:do} & \bkcb & \bkcb \\
 & \R{分度器}{ぶん|ど|き}で \R{角}{かく}を はかれる。 & \pageref{bk:1-20:bundoki} & \bkcb & \bkcb \\
{\sffamily\bfseries\color{bkTheme}1-21} & \R{広}{ひろ}さを、 マスの \R{数}{かず}で くらべられる。 & \pageref{bk:1-21:masu} & \bkcb & \bkcb \\
 & 1\,cm\textsuperscript{2} を もとに、 \R{面積}{めん|せき}を あらわせる。 & \pageref{bk:1-21:cm} & \bkcb & \bkcb \\
 & \R{長方形}{ちょう|ほう|けい}の \R{面積}{めん|せき}を もとめられる。 & \pageref{bk:1-21:kakezan} & \bkcb & \bkcb \\
\end{longtblr}
\endgroup


%% ---------------- つぎのステージへ（第2巻の予告・第5回で追加） ----------------
\bkbacktitle{つぎのステージへ}
この \R{本}{ほん}で、 かぞえる・ はかる・ くらべる ことが できるように なりました。 つぎの 2\R{巻}{かん} 『\R{数}{かず}と \R{量}{りょう}の ことば』では、 こんな ふしぎを しらべます。\par\vspace{3mm}
\newcommand\bkqzcard[4]{% 章, 色, 問い, 絵
  \begin{tikzpicture}[x=1mm, y=1mm]
    \fill[white, draw=bkVol2, line width=1pt, rounded corners=3mm] (0,0) rectangle (66,48);
    \node[fill=bkVol2, text=bkOn2, rounded corners=1mm, font=\sffamily\bfseries\footnotesize, inner xsep=1.6mm, inner ysep=0.8mm] at (9,43) {#1};
    \begin{scope}[shift={(33,24)}] #4 \end{scope}
    \node[anchor=south, text width=60mm, align=center, font=\fontsize{12}{19}\selectfont, text=bkInk] at (33,2.5) {#3};
  \end{tikzpicture}}
\begin{bkcenter}
\bkqzcard{2-2}{}{0より \R{小}{ちい}さい かず って ある？}{%
  \draw[bkInk, line width=0.8pt] (-26,2) -- (26,2);
  \foreach \k in {-5,...,5} {\draw[bkInk, line width=0.6pt] (\k*4.6,0.8) -- (\k*4.6,3.2);}
  \node[font=\sffamily\bfseries\scriptsize] at (0,-1.5) {0};
  \node[font=\sffamily\bfseries\scriptsize, text=bkBlue] at (-13.8,-1.5) {？};
  \node[font=\sffamily\bfseries\scriptsize] at (13.8,-1.5) {3};
  \draw[bkBlue, line width=1pt, -{Stealth[length=2mm]}] (0,6) to[bend right=35] (-13.8,6);}\hfill
\bkqzcard{2-4}{}{ピザを 3\R{人}{にん}で \R{分}{わ}けると？}{%
  \bkpizza{(0,2)}{11}{3}{1}{bkOrange}}\par\vspace{4mm}
\bkqzcard{2-5}{}{「20\%びき」って どれだけ？}{%
  \fill[bkYellow!35, draw=bkStem, line width=0.8pt, rounded corners=1mm] (-18,-6) rectangle (18,10);
  \node[font=\sffamily\bfseries\large, text=bkRed] at (0,4) {20\%};
  \node[font=\sffamily\bfseries\footnotesize, text=bkInk] at (0,-2.5) {わりびき};}\hfill
\bkqzcard{2-11}{}{かずの かわりに \R{文字}{も|じ}？}{%
  \node[font=\fontsize{20}{20}\selectfont, text=bkInk] at (0,2) {□ ＋ 3 ＝ 8};}
\end{bkcenter}
\vspace{3mm}
\bkfukidashi{delta}{1\R{巻}{かん}で のった 7\R{本}{ほん}の \R{路線}{ろ|せん}は、 2\R{巻}{かん}へ つづいて いるよ。 つぎの ステージで まってるね！}

%% ---------------- ふろく（コピーして つかう ページ・第5回で追加） ----------------
%% 見開き（p.236〜237）。 ものさしと 方眼は 実寸（x=1mm）。
\tikzset{bk cut/.style={draw=bkGray, line width=0.5pt, dash pattern=on 1.6pt off 1.3pt}}
\bkbacktitle{ふろく}
{\small コピーして つかう ページです。 コピーは 「100％」で とろう。 \R{点線}{てん|せん}は きる ところ。 はさみは、 おとなと いっしょに つかおう。}\par

\bkuhead{10の わくと おはじき}{（1-2・1-4）}
\begin{bkcenter}
\begin{tikzpicture}[x=1mm, y=1mm]
  \foreach \f in {0,1} {\begin{scope}[shift={(0,-\f*34)}]
    \fill[white] (0,0) rectangle (70,-28);
    \draw[bkInk, line width=0.6pt] (0,-14) -- (70,-14);
    \foreach \x in {14,28,42,56} {\draw[bkInk, line width=0.6pt] (\x,0) -- (\x,-28);}
    \draw[bkInk, line width=1.2pt] (0,0) rectangle (70,-28);
  \end{scope}}
  \begin{scope}[shift={(78,0)}]
    \foreach \i in {0,...,4} {\foreach \j in {0,...,3} {
      \ifnum\j<2 \colorlet{bk@oc}{bkAppleRed}\else\colorlet{bk@oc}{bkBlue}\fi
      \filldraw[fill=bk@oc!30, draw=bk@oc!80!black, line width=0.5pt] (\i*12.4+6.2,-\j*15.5-7.75) circle (4.6);
      \draw[bk cut] (\i*12.4+6.2,-\j*15.5-7.75) circle (5.9);}}
  \end{scope}
\end{tikzpicture}
\end{bkcenter}
{\small \R{上}{うえ}の \R{段}{だん}が いっぱいで 5、 2\R{段}{だん} いっぱいで 10。 おはじきを おいて、 たし\R{算}{ざん}・ひき\R{算}{ざん}を して みよう。}\par

\bkuhead{かずカード}{（1-1〜1-3）}
\begin{bkcenter}
\begin{tikzpicture}[x=1mm, y=1mm]
  \foreach \k in {0,...,11} {\pgfmathtruncatemacro\cx{mod(\k,6)}\pgfmathtruncatemacro\cy{div(\k,6)}
    \begin{scope}[shift={(\cx*23.8,-\cy*29.5)}]
      \fill[white, rounded corners=1.5mm] (0,0) rectangle (21,-27);
      \draw[bk cut, rounded corners=1.5mm] (0,0) rectangle (21,-27);
      \ifnum\k<11
        \node[font=\fontsize{30}{30}\selectfont, text=bkInk] at (10.5,-10.5) {\k};
        \foreach \d in {0,...,9} {\pgfmathtruncatemacro\dx{mod(\d,5)}\pgfmathtruncatemacro\dy{div(\d,5)}
          \ifnum\d<\k \fill[bkVol] (3.1+\dx*3.7,-20.4-\dy*3.4) circle (1.2);
          \else \draw[bkGray!60, line width=0.4pt] (3.1+\dx*3.7,-20.4-\dy*3.4) circle (1.2);\fi}
      \else
        \node[font=\footnotesize, text=bkGray, align=center] at (10.5,-14) {じぶんで\\かこう};
      \fi
    \end{scope}}
\end{tikzpicture}
\end{bkcenter}

\clearpage
\bkuhead{ものさし}{（1-8）}
\begin{bkcenter}
\begin{tikzpicture}[x=1mm, y=1mm]
  \fill[bkYellow!22, rounded corners=1.5mm] (-5,0) rectangle (127,-19);
  \draw[bk cut, rounded corners=1.5mm] (-5,0) rectangle (127,-19);
  \foreach \m in {0,...,120} {\pgfmathsetmacro\bk@tl{ifthenelse(mod(\m,10)==0,6.5,ifthenelse(mod(\m,5)==0,4.6,3))}
    \draw[bkInk, line width=0.35pt] (\m,0) -- (\m,-\bk@tl);}
  \foreach \c in {0,...,12} {\node[font=\sffamily\bfseries\small, text=bkInk] at (\c*10,-10.6) {\c};}
  \node[font=\sffamily\footnotesize, text=bkInk, anchor=east] at (125,-16) {cm};
\end{tikzpicture}
\end{bkcenter}
{\small コピーしたら、 ほんとうの ものさしと ならべて、 1\,cm が おなじか たしかめよう。}\par

\bkuhead{とけい}{（1-10）}
\noindent\begin{minipage}[c]{0.56\linewidth}\centering
\begin{tikzpicture}[x=1mm, y=1mm]
  \fill[white] (0,0) circle (36.5);
  \draw[bk cut] (0,0) circle (36.5);
  \draw[bkInk, line width=1.2pt] (0,0) circle (30);
  \foreach \m in {0,...,59} {\pgfmathsetmacro\bk@a{90-6*\m}
    \ifnum\numexpr\m-5*(\m/5)\relax=0 \draw[bkInk, line width=1pt] (\bk@a:26.6) -- (\bk@a:30);
    \else \draw[bkInk, line width=0.4pt] (\bk@a:28.4) -- (\bk@a:30);\fi}
  \foreach \h in {1,...,12} {\node[font=\sffamily\bfseries\fontsize{13}{13}\selectfont, text=bkInk] at ({90-30*\h}:21.5) {\h};}
  \foreach \h in {1,...,12} {\pgfmathtruncatemacro\bk@mv{5*\h}
    \node[font=\sffamily\bfseries\fontsize{8.5}{9}\selectfont, text=bkBlue!85!black, inner sep=0pt] at ({90-30*\h}:33.4) {\bk@mv};}
  \fill[bkInk] (0,0) circle (1.1);
\end{tikzpicture}
\end{minipage}\hfill
\begin{minipage}[c]{0.42\linewidth}
\begin{tikzpicture}[x=1mm, y=1mm]
  % ながい はり（ふん）
  \begin{scope}[shift={(4,0)}]
    \draw[bk cut, rounded corners=1mm] (-4,-4.4) rectangle (31,4.4);
    \filldraw[fill=bkRed!25, draw=bkRed!85!black, line width=0.8pt, line join=round] (-2.4,-2.2) -- (22,-1.6) -- (27,0) -- (22,1.6) -- (-2.4,2.2) -- cycle;
    \draw[bkRed!85!black, line width=0.6pt] (0,0) circle (0.9);
    \node[anchor=north west, font=\footnotesize, text=bkInk, inner sep=0pt] at (-4,-6) {ながい はり（\R{分}{ふん}）};
  \end{scope}
  % みじかい はり（じ）
  \begin{scope}[shift={(4,-22)}]
    \draw[bk cut, rounded corners=1mm] (-4,-4.8) rectangle (22,4.8);
    \filldraw[fill=bkInk!30, draw=bkInk, line width=0.8pt, line join=round] (-2.4,-2.6) -- (13.5,-2.2) -- (18,0) -- (13.5,2.2) -- (-2.4,2.6) -- cycle;
    \draw[bkInk, line width=0.6pt] (0,0) circle (0.9);
    \node[anchor=north west, font=\footnotesize, text=bkInk, inner sep=0pt] at (-4,-6.4) {みじかい はり（\R{時}{じ}）};
  \end{scope}
\end{tikzpicture}\par\vspace{2mm}
{\small はりを きりとって、 ○の ところに あなを あけ、 とけいの まん\R{中}{なか}に わりピンで とめよう。}
\end{minipage}

\bkuhead{1\,cm の \R{方眼}{ほう|がん}}{（1-19・1-21）}
\begin{bkcenter}
\begin{tikzpicture}[x=1mm, y=1mm]
  \fill[white] (0,0) rectangle (130,-40);
  \draw[bkGray!70, line width=0.4pt] (0,-40) grid[step=10] (130,0);
  \draw[bkInk, line width=0.8pt] (0,0) rectangle (130,-40);
\end{tikzpicture}
\end{bkcenter}
{\small 1マスが 1\,cm\textsuperscript{2}。 \R{形}{かたち}を かいて、 マスを かぞえて \R{広}{ひろ}さを くらべよう。}

%% ---------------- さくいん ----------------
\bkbacktitle{さくいん}
{\small ことばの あとの \R{数}{かず}は、 ページです。}\par\vspace{2mm}
\begin{multicols}{2}\fontsize{12}{20.5}\selectfont\setlength{\parskip}{0pt}\ltjsetruby{stretch=001}% ルビが長い語も、親字を左にそろえる
\directlua{bk.printidx()}
\end{multicols}

%% ---------------- 奥付 ----------------
\bkpadfour
\global\bk@decofalse\bkfolio{false}
\null\vfill
\noindent\begin{minipage}{\linewidth}\fontsize{10}{17}\selectfont
{\sffamily\bfseries\Large\color{bkDeep}\R{物理}{ぶつ|り}を \R{理解}{り|かい}するための \R{数学}{すう|がく}　\R{第}{だい}1\R{巻}{かん}}\par\vspace{1mm}
{\sffamily\bfseries\large かぞえる・はかる}\par\vspace{3mm}
{\color{bkVol}\hrule height 1pt}\vspace{3mm}
\begin{tabular}{@{}l@{\hspace{6mm}}l@{}}
\R{編者}{へん|しゃ} & 「\R{物理}{ぶつ|り}を \R{理解}{り|かい}するための \R{数学}{すう|がく}」\R{編集委員会}{へん|しゅう|い|いん|かい} \\
 & \R{委員長}{い|いん|ちょう}　\R{小学校長}{しょう|がっ|こう|ちょう}（まとめ\R{役}{やく}） \\
 & \R{委員}{い|いん}　\R{数学}{すう|がく}・\R{電気工学}{でん|き|こう|がく}・\R{機械工学}{き|かい|こう|がく}・\R{教育学}{きょう|いく|がく}の \R{大学教授}{だい|がく|きょう|じゅ}、\\
 & \phantom{\R{委員}{い|いん}　}\R{幼稚園教諭}{よう|ち|えん|きょう|ゆ}、 \R{特別支援学校教諭}{とく|べつ|し|えん|がっ|こう|きょう|ゆ}、 \R{心理士}{しん|り|し}、\\
 & \phantom{\R{委員}{い|いん}　}\R{科学解説者}{か|がく|かい|せつ|しゃ}、 \R{小学校教諭}{しょう|がっ|こう|きょう|ゆ} \\
デザイン & ユニバーサルデザイン\R{担当委員}{たん|とう|い|いん} \\
\R{挿絵}{さし|え} & イラスト\R{担当委員}{たん|とう|い|いん} \\
\R{版}{はん} & \R{見本版}{み|ほん|ばん}（\R{第}{だい}5\R{回編集委員会}{かい|へん|しゅう|い|いん|かい}）　2026\R{年}{ねん}10\R{月}{がつ} \\
\end{tabular}\par\vspace{3mm}
{\color{bkVol}\hrule height 1pt}\vspace{2mm}
この \R{本}{ほん}の もんだいの こたえは、 うしろの 「こたえ」に あります。 くわしい とき\R{方}{かた}は、 べっさつの もんだいしゅうと、 \R{先生}{せん|せい}むけの \R{本}{ほん}に のせます。
\end{minipage}
\makeatother
